\begin{abstract}
\noindent We benchmark a set of changes to the density-fitted single-GPU path of GPU4PySCF (a bounded fitted-integral build, a new disk tier next to the existing host-memory tier, cuBLAS contraction fast paths, and a lower-memory analytic gradient) against the unmodified upstream code and against CPU PySCF 2.14.0. The campaign covers \NCasesRun{} cases (SCF energies and gradients from H$_2$O to Taxol with six methods, Hessians, TDA and TDDFT, DF-MP2, CCSD, polarizabilities, shielding, solvation, dispersion, effective-core potentials, geometry optimizations and periodic systems) on P100 (12 and 16\,GB), V100 (16 and 32\,GB) and A100 (80\,GB) GPUs, with two contraction engines (cuTENSOR and the CuPy fallback) and three storage tiers, in \NAuditedRecords{} audited records with accuracy recomputed against an independent CPU result.

\noindent\textbf{Accuracy.} Wherever a result qualified, original and new code agree with the CPU to within the declared gates (worst energy error $8.6\times10^{-9}$~Eh against $10^{-8}$; worst gradient $9.0\times10^{-7}$~Eh/Bohr for the original and $8.5\times10^{-8}$~Eh/Bohr for the new code against $10^{-6}$). The original code fails the gradient gate on Taxol with hybrid functionals; wB97M-V and periodic MgO fail for both versions identically.

\noindent\textbf{Speed.} With cuTENSOR (V100, A100) the new code is as fast as the original (median original/new time ratio \MedianAllAHund{} on the A100 and \MedianAllVHund{} on the V100; per-case ranges \MinAllAHund--\MaxAllAHund{} and \MinAllVHund--\MaxAllVHund), with slowdowns of up to 12\% in individual Hartree--Fock and hybrid-functional cases. Without cuTENSOR (the P100 cards and a forced CuPy engine on V100 and A100) it is faster: 1.2--1.8$\times$ for intermediate-size hybrid-functional jobs and up to \MaxAllAHundCupy$\times$, where the original code pays for \texttt{einsum}. Job types that do not call the modified Fock-build and gradient paths directly (TDDFT, MP2, CCSD, properties, periodic) show no difference, and Hessians and optimizations differ by a few percent at most.

\noindent\textbf{Memory and capacity.} On the V100 the new code uses a median of \MemMedian$\times$ (up to \MemMax$\times$) less tracked GPU memory over the \MemN{} cases whose original peak exceeds 0.3\,GB (medians of 1.8--2.5$\times$ and maxima of 4.0--8.6$\times$ across the six strata). Combined with the host and disk tiers it lets a 12\,GB P100 run valinomycin (1542 basis functions) and an 80\,GB A100 a cluster of 139 water molecules (\BiggestAtoms{} atoms, \BiggestAOs{} basis functions), against a 537-function limit for the original code on the P100 cards (CuPy engine) and 2400 functions on the A100. Larger systems are limited by a gradient tensor of $N_\mathrm{aux}N_\mathrm{occ}^2$ elements, not by the fitted integrals.

\noindent\textbf{Comparison with the CPU.} The GPU is a median of \CPUMedBestAHund$\times$ (A100) and \CPUMedBestVHund$\times$ (V100) faster than the best available CPU lane for cases that take the CPU at least 20\,s, and up to \CPUMaxBestAHund$\times$ (A100) for finite-difference Hessians; a whole-node 48-thread CPU run was found to be slower than a 24-thread, single-BLAS-thread run in 120 of the 123 cases with both lanes (median 2.0$\times$ for the 65 cases above 20\,s on 48 threads, up to 33$\times$ for a Hessian) because of nested BLAS threading, and both are reported.
\end{abstract}
